\documentclass[11pt,a4paper]{article}
\usepackage[T1]{fontenc}
\usepackage{lmodern}
\usepackage[top=2.2cm,bottom=2.4cm,left=2.3cm,right=2.3cm]{geometry}
\usepackage{microtype}
\usepackage{xcolor}
\usepackage{booktabs}
\usepackage{tabularx}
\usepackage{xltabular}   % like tabularx, but the table may continue on the next page
\usepackage{array}
\usepackage{enumitem}
\usepackage[nobottomtitles*]{titlesec}
\usepackage{fancyhdr}
\usepackage[hidelinks]{hyperref}

\definecolor{accent}{HTML}{1F4E79}

% ---- Look -----------------------------------------------------------------------------
\setlength{\parindent}{0pt}
\setlength{\parskip}{0.5em}
\titleformat{\section}{\large\bfseries\color{accent}}{}{0pt}{}[\vspace{-0.5ex}\color{accent}\titlerule]
\titlespacing*{\section}{0pt}{1.4em}{0.6em}
\setlist{nosep,leftmargin=1.4em}
\renewcommand{\arraystretch}{1.25}   % a little air in every table row
% Column types used below: L = left-aligned flexible column that wraps text,
% and C, R = fixed-width columns that are centred or right-aligned.
\newcolumntype{L}{>{\raggedright\arraybackslash}X}
\newcolumntype{C}[1]{>{\centering\arraybackslash}p{#1}}
\newcolumntype{R}[1]{>{\raggedleft\arraybackslash}p{#1}}

\pagestyle{fancy}
\fancyhf{}
\renewcommand{\headrulewidth}{0pt}
\fancyfoot[L]{\footnotesize\color{black!60} STAT 301 \textperiodcentered\ Syllabus, autumn 2026 \textperiodcentered\ subject to change with notice}
\fancyfoot[R]{\footnotesize\color{black!60} \thepage}

\begin{document}

% ---- Title and course facts ---------------------------------------------------------
{\Large\color{accent}\textsc{Example University}\hfill Autumn 2026}\par
{\Huge\bfseries STAT 301: Statistical Learning}\par
\smallskip

% Two columns, label and value. Add a row by copying a line: nothing else changes.
\begin{tabularx}{\textwidth}{@{}>{\bfseries\raggedright\arraybackslash}p{3.9cm}L@{}}
  \toprule
  Instructor      & Prof.\ Alex Morgan, \href{mailto:alex.morgan@example.edu}{alex.morgan@example.edu} \\
  Teaching assistants & Sam Rivera, Chidi Okafor \\
  Lectures        & Mon and Wed, 10:15--11:45, Room 2.04, Hall of Sciences \\
  Tutorial        & Fri, 14:15--15:45, Room 1.12 (choose one of two groups) \\
  Credits         & 4 credits \qquad Language: English \\
  Prerequisites   & Probability (MATH 210) and linear algebra (MATH 120) \\
  Course page     & \href{https://learn.example.edu/stat301}{learn.example.edu/stat301}: slides, problem sets, announcements \\
  \bottomrule
\end{tabularx}

\section{Course description}
How can a computer learn to predict from data, and how do we know whether to trust it?
The course introduces the main methods of supervised learning (linear and logistic
regression, regularisation, trees and ensembles, kernel methods) together with the
statistical ideas that justify them: bias and variance, cross-validation and uncertainty.
Two or three sentences are enough here; say what the course is about and why it matters.

\section{Learning outcomes}
By the end of the course you will be able to:
\begin{enumerate}
  \item explain the bias--variance trade-off and use it to choose between models;
  \item fit, tune and compare regression and classification models in R or Python;
  \item evaluate a model honestly, with held-out data and cross-validation;
  \item derive the estimators of linear and logistic regression and of ridge and lasso;
  \item communicate the result of an analysis to a non-specialist in a short report.
\end{enumerate}
% Write outcomes as things students can do and you can test, starting with a verb.
% Each should be assessed somewhere in the table further down.

\section{Materials}
\textbf{Textbook:} James, Witten, Hastie and Tibshirani, \emph{An Introduction to
Statistical Learning}, 2nd ed.\ (free online version). Further readings are posted on the
course page. Bring a laptop to the tutorials.

\section{Weekly schedule}
% One line per week: Week & Date & Topic & Reading & Due. Dates are the Monday of the
% week. Mark holidays and exams in the Topic column. The table breaks across pages by
% itself and repeats its header.
\begin{xltabular}{\textwidth}{@{}C{1.1cm}C{1.5cm}L>{\raggedright\arraybackslash}p{2.1cm}>{\raggedright\arraybackslash}p{3.0cm}@{}}
  \toprule
  \textbf{Week} & \textbf{Date} & \textbf{Topic} & \textbf{Reading} & \textbf{Due} \\
  \midrule
  \endhead
  \bottomrule
  \endfoot
   1 & Sep 14 & Introduction: prediction, inference, and the statistical learning set-up & ISLR 1--2 & \\
   2 & Sep 21 & Linear regression & ISLR 3 & Problem set 1 \\
   3 & Sep 28 & Classification and logistic regression & ISLR 4 & Problem set 2 \\
   4 & Oct 5  & Resampling: cross-validation and the bootstrap & ISLR 5 & Problem set 3 \\
   5 & Oct 12 & Subset selection and regularisation (ridge, lasso) & ISLR 6 & Project proposal \\
   6 & Oct 19 & Splines and additive models & ISLR 7 & Problem set 4 \\
   7 & Oct 26 & \emph{Midterm exam} (Wed, in class) and review & & Midterm exam \\
   8 & Nov 2  & Decision trees & ISLR 8.1 & \\
   9 & Nov 9  & Bagging, random forests, boosting & ISLR 8.2 & Problem set 5 \\
  10 & Nov 16 & Support vector machines and kernels & ISLR 9 & Problem set 6 \\
  11 & Nov 23 & \emph{No lecture} (mid-term break); work on the project & & \\
  12 & Nov 30 & Unsupervised learning: PCA and clustering & ISLR 12 & Draft report \\
  13 & Dec 7  & Uncertainty, fairness and good practice & Posted notes & Problem set 7 \\
  14 & Dec 14 & Project presentations and course review & & Final report \\
\end{xltabular}

\section{Assessment}
% Weights must add up to 100 %. The last column links each item to the outcomes above.
\begin{tabularx}{\textwidth}{@{}Lp{2.2cm}C{1.6cm}@{}}
  \toprule
  \textbf{Component} & \textbf{Outcomes} & \textbf{Weight} \\
  \midrule
  Problem sets (7, the lowest one is dropped) & 1--4 & 30\,\% \\
  Midterm exam (week 7) & 1, 3, 4 & 25\,\% \\
  Group project: report and presentation & 2, 3, 5 & 30\,\% \\
  Participation in tutorials & 2 & 5\,\% \\
  Final quiz (online, 30 minutes) & 1, 3 & 10\,\% \\
  \midrule
  \textbf{Total} & & \textbf{100\,\%} \\
  \bottomrule
\end{tabularx}

\smallskip
Grades: 90\,\%\,+ is an A, 80\,\% a B, 70\,\% a C, 60\,\% a D. A final grade is
rounded to the nearest whole percent.

\section{Policies}
\textbf{Attendance.} Lectures are not compulsory, tutorials are. You may miss two tutorials
without giving a reason; tell the TAs before the session if you will miss more.

\textbf{Late work.} Each problem set can be handed in up to 48 hours late with a deduction of
10\,\% per day. After that it is not marked, unless you agreed an extension with me before the deadline.

\textbf{Collaboration and academic integrity.} You may discuss problem sets with others, but
you write and code your own solution and name everyone you worked with. Copying, or
submitting work that is not your own, is handled under the university's integrity rules.

\textbf{Use of AI tools.} You may use them to understand a concept or debug code, and you
must say so in your submission. The code and text you hand in must be work you understand
and can explain in an oral check.

\textbf{Accessibility.} If you need accommodations, speak to me in the first two weeks or
contact the Disability Support Office (\href{mailto:access@example.edu}{access@example.edu}).
Everything stays confidential.

\section{Contact and office hours}
\begin{tabularx}{\textwidth}{@{}>{\bfseries\raggedright\arraybackslash}p{3.9cm}>{\raggedright\arraybackslash}p{5.2cm}L@{}}
  \toprule
  Person & \textbf{When} & \textbf{Where / how} \\
  \midrule
  Alex Morgan  & Tue 14:00--16:00 and by appointment & Office 4.12; book on the course page \\
  Sam Rivera   & Thu 10:00--11:00 & Room 1.12 (TA); \href{mailto:s.rivera@example.edu}{s.rivera@example.edu} \\
  Chidi Okafor & Thu 15:00--16:00 & Online (link on the course page) \\
  \bottomrule
\end{tabularx}

\smallskip
For quick questions on the material, use the discussion forum so that everybody sees the
answer; write to me by email for anything personal. I reply within two working days.

\end{document}
